\chapter{Network Codecs and Output Parsing}
\label{ch:codecs}

Every message exchanged between Dynamo's services---routing decisions, token sequences, model outputs, health checks---flows through a single binary codec. Every token generated by the model passes through a parser that extracts structured content (tool calls, reasoning traces, code blocks) from raw text. This chapter examines both: the \texttt{TwoPartCodec} that serializes inter-service communication, and the parser zoo that makes sense of model output. Together, they form Dynamo's data plane---the layer where bytes become meaning. Both are surprisingly rich sources of bugs: the codec's preamble handling contains an integer overflow that bypasses safety checks, and the parser zoo harbors panics, data corruption, and silent data loss across nearly every implementation.

%% ----------------------------------------------------------------
\section{The Wire Format}
\label{sec:codecs:wire}

\marginnote{\emph{Wire format:} the binary layout of data as it appears on the network, before deserialization.}%
Dynamo uses a custom binary protocol for all inter-service TCP communication. The protocol is designed for minimal overhead: no HTTP framing, no JSON parsing on the hot path, and no TLS within the trusted cluster (TLS terminates at the edge).

\begin{whynotbox}{{Why not use Protobuf, gRPC, or JSON?}}
The natural question is: why invent a custom binary protocol when mature alternatives exist? The answer is latency arithmetic. During autoregressive decoding, Dynamo generates one token every few milliseconds per request. Each token must travel from the GPU worker through the codec, across TCP, through the codec on the receiving end, and into the SSE stream to the user. At thousands of concurrent requests, the serialization layer processes millions of messages per second.

\textbf{JSON over HTTP} imposes per-field parsing overhead---field name lookups, string escaping, number parsing---that compounds at this scale. Even fast parsers like \texttt{simd-json} cannot avoid the fundamental cost of scanning for delimiters in a text format. HTTP framing adds further overhead: header parsing, chunked transfer encoding, and connection management.

\textbf{Protocol Buffers / gRPC} are binary and fast, but they impose schema management overhead: \texttt{.proto} files must be kept in sync between sender and receiver, the code generator must run at build time, and the generated code adds compilation cost. gRPC additionally requires HTTP/2 framing, which is overkill for a trusted internal cluster where TLS terminates at the edge.

Dynamo's wire format is simpler than either: a fixed-size 24-byte preamble followed by raw bytes. The receiver reads three integers and immediately knows how many bytes to consume---no parsing, no schema lookup, no framing negotiation. The trade-off is that the format is not self-describing: a version mismatch between sender and receiver causes silent corruption rather than a clean error. For an internal cluster where all services are deployed together, this trade-off is worthwhile.
\end{whynotbox}

Each message consists of three parts:

\begin{enumerate}
  \item \textbf{Preamble} (24 bytes, fixed): three \texttt{u64} fields---the header length, the body length, and an xxHash checksum.
  \item \textbf{Header} (variable length): metadata serialized with MessagePack (a compact binary JSON alternative). Contains routing information, request IDs, and codec version.
  \item \textbf{Body} (variable length): the payload---token sequences, model outputs, or control messages---serialized in a format determined by the message type.
\end{enumerate}

\begin{figure}[h]
\centering
\begin{tikzpicture}[
    field/.style={draw, minimum height=1.0cm, font=\small\ttfamily, anchor=west},
    brace/.style={decorate, decoration={brace, amplitude=6pt, mirror}},
    label/.style={font=\footnotesize},
  ]
  % Preamble fields
  \node[field, fill=blue!12, minimum width=2.8cm] (hlen) at (0,0) {header\_len};
  \node[field, fill=blue!12, minimum width=2.8cm] (blen) at (2.8,0) {body\_len};
  \node[field, fill=blue!12, minimum width=2.8cm] (csum) at (5.6,0) {checksum};

  % Variable-length fields
  \node[field, fill=green!12, minimum width=3.2cm] (header) at (8.4,0) {header data};
  \node[field, fill=orange!12, minimum width=3.8cm] (body) at (11.6,0) {body data};

  % Byte-size labels above
  \node[label, above=2pt of hlen] {8 bytes};
  \node[label, above=2pt of blen] {8 bytes};
  \node[label, above=2pt of csum] {8 bytes};
  \node[label, above=2pt of header] {\texttt{header\_len} bytes};
  \node[label, above=2pt of body] {\texttt{body\_len} bytes};

  % Braces underneath
  \draw[brace] (0,-0.5) -- node[below=8pt, label] {Preamble (24 bytes, fixed)} (8.4,-0.5);
  \draw[brace] (8.4,-0.5) -- node[below=8pt, label] {Payload (variable length)} (15.4,-0.5);

  % Type annotations below braces
  \node[label, text=gray] at (1.4,-1.5) {\texttt{u64} each, big-endian};
  \node[label, text=gray] at (11.9,-1.5) {MessagePack header + raw body};
\end{tikzpicture}
\caption{The \texttt{TwoPartCodec} wire format. The fixed 24-byte preamble contains three \texttt{u64} fields that allow the receiver to determine the exact message size before reading the payload. The checksum is computed only in debug builds; release builds write zero for speed.}
\label{fig:twopart-wire}
\end{figure}

The fixed-size preamble allows the receiver to determine exactly how many bytes to read before parsing begins, enabling zero-allocation reads for small messages that fit in a pre-allocated buffer. This is a length-prefixed framing strategy, common in network protocols from MySQL's packet format to Redis's RESP3. The key advantage over delimiter-based framing (e.g., HTTP's \texttt{\textbackslash r\textbackslash n\textbackslash r\textbackslash n} header terminator) is that the receiver never needs to scan the payload for special characters---it reads exactly the number of bytes declared in the preamble, which means binary payloads require no escaping.

\begin{notebox}
The 24-byte preamble is the first thing the codec reads from the TCP stream. If the preamble is malformed---for example, if the header or body length fields contain values that would overflow when added---the codec must reject the message before allocating memory. This is exactly where Bug~\#3 (Chapter~\ref{ch:findings}) resides: the addition of header length and body length can overflow in release builds, silently bypassing the size check.
\end{notebox}

The wire format is intentionally \emph{not} self-describing: unlike JSON or Protocol Buffers, there are no field names or type tags in the preamble. The header and body are opaque byte sequences whose interpretation depends on the application layer. This keeps the codec layer thin and fast, but it means that a version mismatch between sender and receiver can cause silent data corruption rather than a clean error. Dynamo mitigates this by including a codec version in the MessagePack-encoded header, allowing the receiver to reject messages from incompatible versions.

Dynamo also defines a separate \texttt{TcpRequestMessage} wire format for the request plane, with a more structured layout:

\begin{verbatim}
  endpoint_path_len : u16    (2 bytes)
  endpoint_path     : UTF-8  (variable)
  headers_len       : u16    (2 bytes)
  headers           : JSON   (variable)
  payload_len       : u32    (4 bytes)
  payload           : bytes  (variable)
\end{verbatim}

This format carries routing metadata (the endpoint path) and trace headers inline, enabling the receiving service to dispatch the request without deserializing the payload. The payload itself uses zero-copy slicing via \texttt{bytes::Bytes::slice()}, so the router can forward it to a worker without touching the contents.

The two wire formats serve different roles in Dynamo's architecture. The \texttt{TwoPartCodec} format is used on the \emph{response plane}---the path that carries generated tokens from GPU workers back through the system. Its minimal 24-byte preamble reflects the performance sensitivity of this path: every token generated by every concurrent request passes through this codec. The \texttt{TcpRequestMessage} format is used on the \emph{request plane}---the path that carries user prompts to GPU workers. Requests are less frequent (one per user interaction rather than one per token), so the richer metadata format is affordable.

%% ----------------------------------------------------------------
\section{The TwoPartCodec}
\label{sec:codecs:twopart}

\marginnote{\emph{Tokio codec:} Tokio's \texttt{Decoder}/\texttt{Encoder} traits transform byte streams into typed message streams.}%
With the wire format established, we can now examine the code that reads and writes it. The \texttt{TwoPartCodec} implements Tokio's \texttt{Decoder} and \texttt{Encoder} traits, integrating with Rust's async I/O ecosystem. The core data structure is \texttt{TwoPartMessage}: a pair of \texttt{Bytes} values (header and body) that can be independently empty.

\begin{verbatim}
pub struct TwoPartMessage {
    pub header: Bytes,
    pub data: Bytes,
}
\end{verbatim}

The decoder operates as a state machine with two phases:

\textbf{Phase 1: Read preamble.} The codec checks whether at least 24 bytes are available in the read buffer. If not, it returns \texttt{Ok(None)}, telling Tokio to wait for more data. If 24 bytes are present, it reads three \texttt{u64} values using the \texttt{Buf} trait's \texttt{get\_u64()} method:

\begin{verbatim}
let header_len = cursor.get_u64() as usize;
let body_len   = cursor.get_u64() as usize;
let _checksum  = cursor.get_u64();

let total_len = 24 + header_len + body_len;
\end{verbatim}

\textbf{Phase 2: Read payload.} Using the lengths from the preamble, the codec checks whether \texttt{total\_len} bytes are available. If so, it advances past the 24-byte preamble, then calls \texttt{split\_to()} twice to extract the header and body as separate \texttt{Bytes} values:

\begin{verbatim}
let header = src.split_to(header_len).freeze();
let data   = src.split_to(body_len).freeze();
\end{verbatim}

The \texttt{split\_to()} + \texttt{freeze()} combination is the key to zero-copy operation: it splits the buffer at the given index without copying, then converts the mutable \texttt{BytesMut} slice into an immutable \texttt{Bytes} reference. The critical detail is that \texttt{split\_to(n)} operates in $O(1)$ time---it adjusts internal pointers rather than moving data. The underlying memory is shared between the original buffer and the split-off piece via reference counting.

\marginnote{\emph{xxHash:} a non-cryptographic hash function optimized for speed, producing 64-bit hashes at $\sim$30\,GB/s on modern CPUs.}%
The choice of xxHash-64 (\texttt{xxh3\_64} from the \texttt{xxhash-rust} crate) for the checksum is deliberate: it is one of the fastest non-cryptographic hash functions available, processing data at memory bandwidth on modern CPUs. Since the checksum is only computed in debug builds, its purpose is catching bugs during development, not protecting against adversarial corruption---TCP's own 16-bit checksum handles the latter (imperfectly, but adequately for a trusted cluster network).

The encoder reverses this process: given a \texttt{TwoPartMessage}, it writes the three \texttt{u64} preamble fields followed by the header and body bytes. Checksum computation is conditional: in debug builds, it computes an xxHash-64 over the concatenation of header and body; in release builds, it writes a zero checksum for speed, relying on TCP's own checksumming for integrity.

The \texttt{TwoPartMessage} type also provides convenience constructors for common patterns: \texttt{from\_header()} for control messages with no body, \texttt{from\_data()} for payloads with no routing header, and \texttt{into\_message\_type()} which classifies the message into one of four variants (\texttt{HeaderOnly}, \texttt{DataOnly}, \texttt{HeaderAndData}, \texttt{Empty}) for pattern matching in downstream handlers.

\begin{notebox}
The codec supports multiple messages concatenated in a single TCP buffer. Tokio's \texttt{FramedRead} adapter calls \texttt{decode()} in a loop until it returns \texttt{Ok(None)}, naturally handling message pipelining. This is important for performance: TCP's Nagle algorithm may coalesce multiple small messages into a single network packet, and the codec must correctly split them apart.
\end{notebox}

\begin{bugbox}
\textbf{Bug \#3: Integer overflow in preamble handling.} The codec reads \texttt{header\_len} and \texttt{body\_len} as \texttt{u64} values from the network, then computes \texttt{let total\_len = 24 + header\_len + body\_len}. All three operands are \texttt{usize} (64-bit on modern systems), but the addition can still overflow. Consider \texttt{header\_len = u64::MAX} and \texttt{body\_len = 1}: their sum wraps to zero, so \texttt{total\_len} becomes 24. If \texttt{max\_message\_size} is set to, say, 16\,MB, the check \texttt{total\_len > max\_size} passes (24 < 16\,MB), and the codec proceeds to read from the buffer using the original, enormous \texttt{header\_len}---causing an out-of-bounds read or an allocation of $2^{64} - 1$ bytes.

In Rust's debug build, the addition panics (denial of service). In release builds, the addition wraps silently via two's complement, potentially allowing an attacker-controlled payload to bypass length validation entirely. The fix is to use \texttt{checked\_add()}, which returns \texttt{None} on overflow:

\begin{verbatim}
let total_len = 24_usize
    .checked_add(header_len)
    .and_then(|n| n.checked_add(body_len))
    .ok_or(TwoPartCodecError::InvalidMessage(
        "Length overflow".to_string(),
    ))?;
\end{verbatim}

Note that the debug-mode checksum verification path \emph{does} use \texttt{checked\_add}---the developers were aware of the overflow risk in that code path but missed the earlier, more critical one in the main decode logic. This is a common pattern in security bugs: the fix is applied in one place but not another.
\end{bugbox}

%% ----------------------------------------------------------------
\section{Zero-Copy Decoding}
\label{sec:codecs:zerocopy}

\marginnote{\emph{Zero-copy:} processing data without copying it to a new memory location, operating directly on the original buffer.}%
The previous section described \emph{what} the codec reads; this section explains \emph{how} it avoids the most expensive part of reading: copying data. For high-throughput message processing, copying data between buffers is a significant cost. A single \texttt{memcpy} of a 4\,KB token sequence takes roughly 200 nanoseconds---trivial in isolation, but at one million messages per second, that is 200 milliseconds of CPU time per second consumed purely by copying. The codec is designed to eliminate this cost.

The mechanism is Rust's \texttt{bytes} crate, which provides two core types:

\begin{itemize}
  \item \texttt{BytesMut}: a mutable, growable byte buffer with reference-counted backing storage. Tokio's framed reader fills this buffer with data from the TCP socket.
  \item \texttt{Bytes}: an immutable, cheaply cloneable view into a \texttt{BytesMut} buffer. Cloning a \texttt{Bytes} value increments a reference count---it does not copy the underlying data.
\end{itemize}

The \texttt{Buf} and \texttt{BufMut} traits provide cursor-like access to these buffers. \texttt{Buf::get\_u64()} reads 8 bytes and advances the cursor; \texttt{BufMut::put\_u64()} writes 8 bytes and advances. These traits allow the codec to read and write structured fields without manual index arithmetic.

The zero-copy flow works as follows:

\begin{enumerate}
  \item Tokio reads TCP data into a \texttt{BytesMut} buffer (a single allocation, reused across messages).
  \item The codec's \texttt{decode()} method calls \texttt{src.split\_to(header\_len)}, which cleaves the buffer at the given offset. The original buffer retains bytes after the split point; the returned \texttt{BytesMut} owns the bytes before it. No data is copied.
  \item The codec calls \texttt{.freeze()} on the split result, converting \texttt{BytesMut} to \texttt{Bytes}. This makes the buffer immutable and enables cheap cloning.
  \item The resulting \texttt{Bytes} header and body are passed to the deserialization layer, which can further \texttt{.slice()} them without copying.
\end{enumerate}

The trade-off is that zero-copy decoding holds a reference to the original buffer's backing allocation, preventing it from being reused until all slices are dropped. For long-lived messages (e.g., streaming responses that persist for the duration of generation), this can increase peak memory usage. For the response plane, Dynamo also provides a \texttt{ZeroCopyTcpDecoder} that uses the same \texttt{split\_to/freeze} pattern for the request message format, extracting the payload as \texttt{Bytes::slice()} directly from the read buffer.

The response-plane codec (\texttt{TcpResponseCodec}) uses an even simpler format: a 4-byte \texttt{u32} length prefix followed by the data bytes. This minimalism reflects the asymmetry of inference serving: requests carry rich metadata (endpoint paths, trace headers, routing information), while responses are often just serialized token sequences that need minimal framing. The response codec's \texttt{decode()} method follows the same two-phase pattern as \texttt{TwoPartCodec}: peek at the length, check availability, then \texttt{split\_to().freeze()}.

\begin{notebox}
The \texttt{TcpRequestMessage} format uses \texttt{u16} for endpoint path and header lengths, limiting each to 65\,535 bytes, and \texttt{u32} for payload length, limiting payloads to 4\,GB. These limits are enforced both in encoding (reject oversized inputs before writing) and decoding (reject malformed length fields before allocating). This belt-and-suspenders approach prevents both memory exhaustion from malicious inputs and silent truncation from overflow.
\end{notebox}

\begin{insightbox}
The zero-copy pattern explains why Dynamo uses \texttt{Bytes} pervasively rather than \texttt{Vec<u8>}. A \texttt{Vec<u8>} must be cloned with a full copy; a \texttt{Bytes} is cloned with an atomic reference count increment (a single CPU instruction on x86). When a single incoming message is fanned out to multiple downstream consumers---as happens during broadcast or multi-worker routing---the difference is substantial.
\end{insightbox}

%% ----------------------------------------------------------------
\section{TCP Connection Management}
\label{sec:codecs:tcp}

\marginnote{\emph{Framed adapter:} Tokio's \texttt{Framed\-Read}/\texttt{Framed\-Write} wrappers apply a codec to a raw byte stream, producing a typed \texttt{Stream} of messages.}%
The codec knows how to encode and decode individual messages, but it does not manage the connections those messages flow over. That is the responsibility of the TCP connection layer. Dynamo maintains persistent TCP connections between services, using connection pooling to amortize the cost of TCP handshakes. Each connection is wrapped in a Tokio \texttt{Framed} adapter that applies the \texttt{TwoPartCodec} to the raw byte stream, producing a stream of typed \texttt{TwoPartMessage} values.

The server side (\texttt{TcpStreamServer}) binds to a port using the \texttt{socket2} crate for fine-grained control over socket options (e.g., \texttt{SO\_REUSEADDR}). When a client connects, the server splits the TCP stream into read and write halves, wraps each in a \texttt{FramedRead} or \texttt{FramedWrite} with the appropriate codec, and hands them to the pipeline.

Connection establishment uses a ``call home'' handshake: the worker connects to the router (rather than the router connecting to the worker), sends a \texttt{CallHomeHandshake} identifying itself, and the router registers the connection. This design simplifies firewall rules in containerized deployments---workers initiate outbound connections, which is generally permitted by default.

Connection health is monitored through heartbeat messages: the router periodically sends a ping, and the worker must respond within a timeout. If a worker fails to respond, the router marks it as unhealthy and stops routing requests to it, then notifies the service discovery layer (etcd) to trigger a replacement.

Backpressure is handled at the connection level: if a worker's incoming message queue is full, the TCP flow control mechanism (window size) naturally slows the sender. The codec does not implement application-level backpressure beyond what TCP provides, relying on the OS network stack to regulate flow.

\begin{insightbox}
The ``call home'' design inverts the usual client-server relationship: the worker is the TCP client even though it is the service provider. This is a pragmatic choice for Kubernetes deployments where workers are ephemeral pods with dynamic IP addresses. Rather than requiring the router to discover worker addresses (via DNS, service mesh, or etcd polling), the worker connects to the router's well-known address at startup. The handshake message carries the worker's capabilities (GPU type, available memory, supported models), allowing the router to begin routing immediately.
\end{insightbox}

The Tokio \texttt{FramedRead} and \texttt{FramedWrite} adapters deserve special mention. They bridge the gap between Tokio's byte-oriented \texttt{AsyncRead}/\texttt{AsyncWrite} traits and the message-oriented \texttt{Stream}/\texttt{Sink} traits. The \texttt{FramedRead} adapter maintains an internal \texttt{BytesMut} buffer, reads bytes from the socket into it, then repeatedly calls the codec's \texttt{decode()} method until it returns \texttt{Ok(None)} (not enough data) or an error. This loop handles the case where a single TCP read delivers multiple complete messages---common under high load when the OS batches socket reads.

%% ----------------------------------------------------------------
\section{The Parser Zoo}
\label{sec:codecs:parsers}

\marginnote{\emph{Parser zoo:} Dynamo's collection of 17+ model-specific output parsers, each handling a different model family's output format.}%
The codec and TCP layers handle \emph{transport}: getting bytes reliably from one service to another. But once a message arrives, a second challenge begins: making sense of the \emph{content}. For model output in particular, this is where things get messy.

Large language models do not produce structured output. They produce raw text---a stream of tokens decoded back into characters. But applications need structure: tool call invocations with JSON arguments, reasoning traces separated from final answers, code blocks with language annotations. The \emph{parser zoo} is Dynamo's collection of model-specific parsers that extract this structure from raw text.

Why are there so many parsers? Because there is no standard tool-call format. The OpenAI API defines a \emph{response} format for tool calls (a JSON object with \texttt{name} and \texttt{arguments} fields), but it does not standardize the \emph{output} format---the raw text that the model generates before it is parsed into that response structure. Each model family invented its own convention for how the model should emit tool calls in its token stream:

\begin{itemize}
  \item \textbf{Hermes/Llama3/Mistral/Phi4}: JSON-based formats with start/end delimiter tokens (e.g., \texttt{<TOOLCALL>...\allowbreak</TOOLCALL>}).
  \item \textbf{DeepSeek V3}: custom delimiters with Unicode special characters: \texttt{<|tool\_call\_begin|>}.
  \item \textbf{DeepSeek V3.1/V3.2 (DSML)}: ``DeepSeek Markup Language,'' a full XML-like format with Unicode delimiters and typed parameters.
  \item \textbf{GLM-4.7} (Zhipu AI): XML with nested tags such as \texttt{<tool\_call>}, \texttt{<arg\_key>}, and \texttt{<arg\_value>}.
  \item \textbf{Kimi K2} (Moonshot AI): section-delimited format with \texttt{<|tool\_call\_begin|>} tokens and \texttt{functions.name:index} identifiers.
  \item \textbf{XML/Qwen3 Coder}: generic XML with \texttt{<function=name>\allowbreak <parameter=key>value\allowbreak </parameter>\allowbreak </function>}.
  \item \textbf{Pythonic}: Python function call syntax---\texttt{[get\_weather(city="SF")]}.
  \item \textbf{Harmony}: an async JSON-based format requiring concurrent processing.
  \item \textbf{MiniMax M2}: yet another JSON variant with model-specific quirks.
\end{itemize}

This diversity is the price of a decentralized ecosystem: without a single authority mandating a wire format, each model trainer chose whatever format their training data and chat templates made convenient. The serving infrastructure must handle them all.

\begin{whynotbox}{Why so many parser implementations instead of one unified parser?}
Given the bug density of the parser zoo (the majority of bugs cataloged in Chapter~\ref{ch:findings} reside in parsers), a natural design question is: why not define a single, well-tested parser that all models funnel through?

The answer is that the formats are \emph{genuinely different}, not superficial variations. Hermes uses JSON between XML-like delimiters. DeepSeek DSML uses a full XML-like markup with typed attributes. GLM-4.7 uses key-value tag pairs. The Pythonic format is literal Python syntax. You cannot write a single regex or grammar that handles all of these---the syntactic structures are incompatible. A ``unified parser'' would need to be a full parser combinator framework with pluggable grammars, which is effectively what the current enum-dispatched architecture already is, just with hand-written parsers instead of generated ones.

The deeper problem is that model trainers control the output format, not the serving infrastructure. A model is trained to emit tool calls in a specific syntax baked into its training data and chat template. Changing the format would require retraining the model. The serving layer has no choice but to meet each model where it is.

What \emph{could} be unified---and would reduce bugs---is the streaming state machine logic. Every parser faces the same streaming challenge (detecting partial delimiters, buffering across token boundaries, handling boundary mismatches), and each reimplements it independently. A shared streaming harness with pluggable delimiter sets would eliminate an entire class of bugs.
\end{whynotbox}

The parser architecture unifies these behind two trait interfaces:

\textbf{Tool call parsing} uses a dispatcher pattern. A static \texttt{PARSER\_MAP} (initialized via \texttt{OnceLock}) maps parser names to \texttt{ToolCallConfig} structs. Each config specifies the parser variant via a \texttt{ParserConfig} enum:

\begin{verbatim}
pub enum ParserConfig {
    Json(JsonParserConfig),
    Pythonic,
    Xml(XmlParserConfig),
    Dsml(DsmlParserConfig),
    Glm47(Glm47ParserConfig),
    KimiK2(KimiK2ParserConfig),
    Harmony(JsonParserConfig),
    Typescript,  // Not yet implemented
}
\end{verbatim}

The entry point, \texttt{try\_tool\_call\_parse()}, dispatches on this enum to the appropriate parser implementation. Each parser returns a \texttt{(Vec<ToolCallResponse>, Option<String>)} tuple: the parsed tool calls and any remaining normal text. A \texttt{ToolCallResponse} contains three fields: a unique ID, a type (currently always \texttt{Function}), and a \texttt{CalledFunction} with the function name and JSON-serialized arguments:

\begin{verbatim}
pub struct ToolCallResponse {
    pub id: String,
    pub tp: ToolCallType,
    pub function: CalledFunction,
}
pub struct CalledFunction {
    pub name: String,
    pub arguments: String,  // JSON
}
\end{verbatim}

\textbf{Reasoning parsing} uses a trait-object pattern. The \texttt{ReasoningParser} trait defines two methods:

\begin{verbatim}
pub trait ReasoningParser: Send + Debug {
    fn detect_and_parse_reasoning(
        &mut self, text: &str, token_ids: &[u32]
    ) -> ParserResult;

    fn parse_reasoning_streaming_incremental(
        &mut self, text: &str, token_ids: &[u32]
    ) -> ParserResult;
}
\end{verbatim}

The first method handles one-shot (complete) inputs; the second handles streaming, returning only the \emph{delta} (newly discovered text) for each chunk. The \texttt{ParserResult} struct splits text into \texttt{reasoning\_text} and \texttt{normal\_text}. Implementations include \texttt{BasicReasoningParser} (handles \texttt{<think>...</think>} with configurable delimiters, used by DeepSeek R1, Qwen3, Kimi K2.5, and Mistral), \texttt{GraniteReasoningParser} (natural-language delimiters like ``Here's my thought process:''), \texttt{GptOssReasoningParser} (GPT-compatible reasoning extraction), and \texttt{MiniMaxAppendThinkParser}. The diversity of delimiter styles is remarkable: DeepSeek uses XML tags, Kimi (original) uses Unicode triangular brackets (\texttt{<<think>>}), Mistral uses bracketed tags (\texttt{[THINK]...[/THINK]}), and Granite uses natural English sentences.

The \texttt{ReasoningParser} trait is wrapped in a \texttt{ReasoningParserWrapper} that holds a \texttt{Box<dyn ReasoningParser>}, enabling dynamic dispatch. A factory method \texttt{get\_reasoning\_parser\_from\_name()} looks up the parser type by string name from a global \texttt{OnceLock}-backed map and returns the appropriate wrapper. Unknown names fall back to the basic \texttt{<think>...</think>} parser with a warning---a resilient default that prevents misconfiguration from crashing the system.

\begin{notebox}
The parser zoo currently maps 17 named parsers to 9 distinct implementations. Several models share implementations: Nemotron Nano reuses Qwen3 Coder's parser, GLM-4.5 reuses the basic \texttt{<think>} parser. The \texttt{OnceLock}-based global map means these mappings are established once and cached for the process lifetime---which becomes a source of bugs when the mapping depends on runtime configuration (see Bug~\#11, the Kimi K2 \texttt{OnceLock} caching issue).
\end{notebox}

%% ----------------------------------------------------------------
\section{Tool Call Parsing}
\label{sec:codecs:toolcalls}

\marginnote{\emph{Tool call:} a structured invocation embedded in model output, specifying a function name and JSON arguments.}%
Tool calls are the primary structured output that parsers must extract. When a model is prompted with a list of available tools, it may respond with a tool invocation---a function name and arguments---embedded in its text output. This section walks through how the major parsers handle this, and where they break.

Each parser implements three functions that mirror the lifecycle of a streaming tool call:

\begin{enumerate}
  \item \texttt{detect\_tool\_call\_start(chunk)} --- returns \texttt{true} if the chunk contains (or ends with a prefix of) the tool call start token. This is used by the streaming layer to switch from ``emit text'' mode to ``buffer tool call'' mode.
  \item \texttt{find\_tool\_call\_end\_position(chunk)} --- returns the byte offset after the end token, or the chunk length if the end token has not yet appeared. This tells the streaming layer how much of the buffered text constitutes a complete tool call that can be handed to the parser.
  \item \texttt{try\_tool\_call\_parse(message)} --- the main parsing function, called once a complete tool call block has been accumulated. Returns structured tool calls and any non-tool-call text that was adjacent to the tool call block.
\end{enumerate}

The partial-prefix detection in step 1 deserves attention. Consider a streaming scenario where the model outputs the text ``Some text <tool\_c'' in one chunk. The start token \texttt{<tool\_call>} is not fully present, but the chunk \emph{ends with a prefix of it}. The parser must recognize this partial match and begin buffering, because the next chunk might complete the token. All parsers use a shared utility function \texttt{chunk\_ends\_with\_token\_prefix()} for this check, which iterates over successively longer suffixes of the chunk to see if any is a prefix of the start token. This function was itself a source of bugs before being fixed: the original byte-based implementation panicked on multi-byte UTF-8 start tokens (e.g., the Chinese CJK characters meaning ``tool'' characters used in some GLM configurations).

\subsection{The XML Parser Family}

The XML parser handles formats like Qwen3 Coder, where tool calls look like:

\begin{verbatim}
<tool_call>
<function=get_weather>
<parameter=city>San Francisco</parameter>
<parameter=unit>fahrenheit</parameter>
</function>
</tool_call>
\end{verbatim}

The parser uses regex to extract function blocks and parameter blocks from within each \texttt{<tool\_call>...</tool\_call>} region. Parameters are typed using the tool's JSON Schema: if the schema says \texttt{city} is a string and \texttt{count} is an integer, the parser coerces accordingly. Without a schema, all values are returned as strings. The type coercion logic is extensive, recognizing aliases like \texttt{int}/\texttt{integer}/\texttt{int32}/\texttt{int64}, \texttt{float}/\texttt{number}/\texttt{double}, \texttt{boolean}/\texttt{bool}/\texttt{binary}, \texttt{object}/\texttt{dict}/\texttt{dictionary}, and \texttt{array}/\texttt{arr}/\texttt{list}.

The parser is also deliberately fault-tolerant with missing closing tags. If \texttt{</parameter>} is absent, the regex's non-greedy match extends to the next \texttt{<parameter=} or \texttt{</function>} tag, capturing the value with trailing markup. If both \texttt{</parameter>} and \texttt{</function>} are missing, the value extends to the \texttt{</tool\_call>} tag---including the closing tag itself in the value. This matches the behavior of the upstream SGLang Python implementation, trading correctness for resilience against malformed model output.

The XML parser contains two bugs found by fuzzing:

\begin{bugbox}
\textbf{Bug \#13: \texttt{strip\_quotes} panics on single-character input.} The utility function \texttt{strip\_quotes()} removes surrounding quotes from a string. It checks \texttt{starts\_with('"') \&\& ends\_with('"')}, then returns \texttt{\&trimmed[1..trimmed.len()-1]}. But when the input is a single quote character \texttt{"}, both conditions are true (a single-char string both starts and ends with that character), and the slice becomes \texttt{[1..0]}---an invalid range. Rust panics with ``begin <= end'' violated. This is a classic off-by-one: the code assumes quoted strings are at least two characters long.
\end{bugbox}

\begin{bugbox}
\textbf{Bug \#12: \texttt{try\_literal\_eval} corrupts string values.} The XML parser falls back to a Python \texttt{ast.literal\_eval}-style parser for non-JSON values. This function performs global string replacements: \texttt{.replace("True", "true")}, \texttt{.replace("False", "false")}, \texttt{.replace("None", "null")}. These replacements are not context-aware---they corrupt substrings within quoted values. The input \texttt{\{'location': 'TrueNorth'\}} becomes \texttt{\{"location": "trueNorth"\}}, silently corrupting the string ``TrueNorth'' to ``trueNorth''. Similarly, ``Falsehood'' becomes ``falsehood'' and ``NoneAvailable'' becomes ``nullAvailable''.
\end{bugbox}

\subsection{The GLM-4.7 Parser}

GLM-4.7 (Zhipu AI) uses a distinct XML dialect where arguments are structured as key-value tag pairs rather than attributes:

\begin{verbatim}
<tool_call>get_weather
  <arg_key>location</arg_key>
  <arg_value>San Francisco</arg_value>
</tool_call>
\end{verbatim}

The parser extracts the function name as everything between \texttt{<tool\_call>} and the first \texttt{<arg\_key>}, then uses a regex to match \texttt{<arg\_key>...</arg\_key><arg\_value>...</arg\_value>} pairs. Values are coerced using the tool's schema (integers, booleans, arrays, etc.), with XML entity decoding (\texttt{\&lt;} $\rightarrow$ \texttt{<}) applied before parsing. The parser also handles multi-line argument values, which is essential for code-generation tool calls.

\begin{bugbox}
\textbf{Bug \#1: UTF-8 boundary panic in GLM-4.7 parser.} The parser extracts the function name by trimming whitespace from the content between \texttt{<tool\_call>} and the first \texttt{<arg\_key>}. It then computes the arguments section as \texttt{content[function\_name.len()..]}\,---using the \emph{trimmed} name's byte length as an offset into the \emph{untrimmed} content. When there is leading whitespace and the function name contains multi-byte UTF-8 characters (e.g., CJK characters like CJK characters meaning ``get weather''), the byte offset lands in the middle of a UTF-8 code point, and Rust panics with ``byte index is not a char boundary.'' For ASCII function names with leading whitespace, the offset is merely wrong (off by the whitespace length), but the regex still finds arguments---silently producing correct results by accident. With multi-byte characters, the bug becomes a crash.
\end{bugbox}

\subsection{The Pythonic Parser}

The Pythonic parser handles models (e.g., Jamba, some Llama variants) that emit tool calls as Python expressions:

\begin{verbatim}
[get_weather(city="San Francisco"), set_alarm(time="8:00")]
\end{verbatim}

Rather than writing a custom parser, Dynamo uses the \texttt{rustpython-parser} crate to parse this as a Python AST. The expression is parsed in \texttt{Mode::Expression}, and the AST is walked to extract function names and keyword arguments. This is elegant---it handles nested data structures (lists, dicts) and Python literals (booleans, None) natively. The AST walker recursively converts Python constants to JSON values: \texttt{True} $\rightarrow$ \texttt{true}, \texttt{None} $\rightarrow$ \texttt{null}, Python integers to JSON numbers (with fallback to strings for values outside the \texttt{i64}/\texttt{u64} range), and Python dicts and lists to their JSON equivalents.

A pre-processing regex identifies candidate tool call blocks in the input. The pattern matches the outer \texttt{[...]} list containing one or more \texttt{func(key=val)} calls. The regex is compiled once via \texttt{OnceLock} and reused across invocations. Text before the first match is returned as ``normal text''; the matched block is passed to the Python parser.

However, the approach has bugs:

\begin{bugbox}
\textbf{Bug \#7: Pythonic parser prefix leak and text loss.} The parser has multiple issues. First, it extracts ``normal text'' as \texttt{stripped.split(\&matches[0]).next()}---which returns only text \emph{before} the first match. Any text after the tool call is silently dropped. Second, adjacent identifier characters are absorbed into the function name, corrupting the extracted name. Third, the regex pattern requires at least one \texttt{key=value} argument, so a zero-argument call like \texttt{[get\_time()]} is treated as normal text rather than a tool call.
\end{bugbox}

\subsection{The JSON Parser Family}

The JSON-based parsers handle the largest number of model families: Hermes, Llama 3, Mistral, Phi4, Nemotron, Jamba, MiniMax M2, and the DeepSeek V3/V3.1 variants. The common pattern is: find text between configurable start and end delimiter tokens, then parse the enclosed content as JSON. The JSON object is expected to contain a function name key and an arguments key:

\begin{verbatim}
<TOOLCALL>{"name": "get_weather",
           "arguments": {"city": "SF"}}</TOOLCALL>
\end{verbatim}

DeepSeek V3 uses a more complex format with Unicode delimiters and a type/separator structure:

\begin{verbatim}
<|tool_call_begin|>function<|tool_sep|>get_weather
```json
{"city": "SF"}
```<|tool_call_end|>
\end{verbatim}

The JSON parser extracts blocks using regex, then calls \texttt{serde\_json::from\_str} to parse each block. DeepSeek V3.1 adds a further complication: it uses both ``individual call'' delimiters (\texttt{tool\_call\_begin/end}) and ``wrapper'' delimiters (\texttt{tool\_calls\_begin/end}) that enclose multiple calls. The parser must filter delimiter tokens to find the right level of nesting. The \texttt{JsonParserConfig} is accordingly flexible, with separate fields for start tokens, end tokens, separator tokens, function name keys, and argument keys---all as \texttt{Vec<String>} to accommodate models that use different keys (e.g., \texttt{"name"} vs.\ \texttt{"function"}, \texttt{"arguments"} vs.\ \texttt{"parameters"}).

Note that DeepSeek V3's coverage from generic fuzz targets was zero---the Unicode delimiter tokens are unlikely to appear in random input. Achieving meaningful coverage required constructing parser-specific fuzz targets with structured inputs that include the correct delimiter tokens, a challenge we discuss further in Chapter~\ref{ch:fuzzing}.

\subsection{The DSML Parser}

DeepSeek V3.2 introduced DSML (DeepSeek Markup Language), a full XML-like format with Unicode-heavy delimiters:

\begin{verbatim}
<|DSML|function_calls>
<|DSML|invoke name="get_weather">
<|DSML|parameter name="city" string="true">
  San Francisco
</|DSML|parameter>
</|DSML|invoke>
</|DSML|function_calls>
\end{verbatim}

Each parameter carries a \texttt{string="true|false"} attribute indicating whether the value should be treated as a literal string or parsed as JSON. The parser uses regex to extract invoke blocks and parameter blocks, then conditionally parses values based on this flag. The DSML format is notable for its use of fullwidth Unicode pipe characters (\texttt{|}, U+FF5C) in its delimiters---these serve as visual markers that are extremely unlikely to appear in natural text, reducing false-positive detection. The downside is that these characters are multi-byte in UTF-8 (3 bytes each), making byte-level string manipulation error-prone.

\begin{bugbox}
\textbf{Bug \#10: DSML parser silently drops parameters.} The regex that matches \texttt{<|DSML|parameter name="..."} requires the \texttt{string} attribute to be present. If the model omits the \texttt{string} attribute (which happens for non-string parameters in some model configurations), the regex does not match, and the parameter is silently dropped from the parsed output. The tool call succeeds but with missing arguments---a silent data loss bug.
\end{bugbox}

\subsection{The Kimi K2 Parser}

Kimi K2 (Moonshot AI) uses section delimiters and a \texttt{functions.name:index} naming convention:

\begin{verbatim}
<|tool_calls_section_begin|>
<|tool_call_begin|>
functions.get_weather:0
<|tool_call_argument_begin|>
{"city": "SF"}
<|tool_call_end|>
<|tool_calls_section_end|>
\end{verbatim}

The function ID is parsed with a secondary regex that strips the optional \texttt{functions.} prefix and extracts the name and index components. The arguments are expected to be a JSON object between the \texttt{argument\_begin} and \texttt{call\_end} tokens.

The parser compiles its main regex from the config tokens using \texttt{OnceLock}, caching it for the process lifetime. The pattern captures two named groups: \texttt{function\_id} (matching the \texttt{[\textbackslash w.]+:\textbackslash d+} format) and \texttt{arguments} (the JSON object). This leads to:

\begin{bugbox}
\textbf{Bug \#11: \texttt{OnceLock} caches stale regex.} The Kimi K2 parser stores its compiled regex in a \texttt{static OnceLock<Regex>}. The regex pattern is built from the \texttt{KimiK2ParserConfig}'s delimiter tokens. But \texttt{OnceLock::get\_or\_init} only runs the initialization closure once---the first time it is called. If the process handles requests for multiple Kimi K2 model variants with different delimiter configurations (e.g., singular \texttt{tool\_call\_section\_begin} vs.\ plural \texttt{tool\_calls\_section\_begin}), the regex cached from the first variant is silently used for all subsequent variants. Tool calls from the second variant fail to match.
\end{bugbox}

\subsection{The Granite Reasoning Parser}

\marginnote{\emph{Streaming mismatch:} when one-shot and streaming parsing produce different results for the same input.}%
The Granite parser (IBM's model family) separates reasoning from content using natural-language delimiters: ``Here's my thought process:'' and ``Here's my response:''. The parser accepts multiple variants of each delimiter (e.g., ``Here's my thought process:'' and ``Here is my thought process:'') to handle the model's non-deterministic phrasing. Unlike XML tags, these delimiters can appear as substrings of normal text, making boundary detection fragile. The parser maintains internal state (\texttt{in\_reasoning}, \texttt{stripped\_think\_start}, a text buffer) to track its position in the reasoning/content split across streaming chunks.

\begin{bugbox}
\textbf{Bug \#8: Granite streaming mismatches.} The one-shot parser (\texttt{detect\_and\_parse\_reasoning}) and the streaming parser (\texttt{parse\_reasoning\_streaming\_incremental}) produce different results for the same input. The one-shot parser uses \texttt{.find()} to locate delimiters in the complete text; the streaming parser accumulates text in a buffer and checks for delimiters at each token boundary. When a delimiter spans a token boundary (e.g., ``Here's my '' in one token and ``thought process:'' in the next), the streaming parser may detect it at a different point than the one-shot parser, causing the split between reasoning and content to differ.
\end{bugbox}

The MiniMax Append Think parser exhibits a similar class of streaming mismatch (Bug~\#9): its one-shot and streaming modes disagree on where reasoning text ends and normal text begins, particularly when the \texttt{<think>} block is empty or contains only whitespace.

\begin{insightbox}
The streaming parsing problem is fundamentally harder than batch parsing. In batch mode, the parser sees the entire input and can use global context to locate boundaries. In streaming mode, it sees one token at a time and must decide---before the next token arrives---whether the current token is inside a tool call, outside one, or at a boundary. A single misclassified token can cause the parser to lose synchronization: buffering text that is not a tool call (causing latency), or emitting raw text that was meant to be a structured argument (causing data corruption). Every streaming parser bug we found involves these boundary transitions. This is not a coincidence: streaming parsers are state machines operating on partial information, and state machines are precisely the class of program where fuzzers excel at finding unexpected state transitions.
\end{insightbox}

%% ----------------------------------------------------------------
\section{Summary}
\label{sec:codecs:summary}

\begin{itemize}
  \item The \textbf{wire format} uses a 24-byte fixed preamble (header length, body length, checksum) followed by variable-length header and body payloads. A separate request-plane format adds endpoint routing and trace headers.
  \item The \textbf{TwoPartCodec} implements Tokio's codec traits for async encoding and decoding, with a two-phase state machine (preamble, then payload). The checksum is only computed in debug builds; release builds write zero for performance.
  \item \textbf{Zero-copy decoding} via \texttt{BytesMut::split\_to().freeze()} avoids allocation on the hot path, using reference-counted byte slices that can be cheaply cloned and sliced without copying data.
  \item \textbf{TCP connection management} uses persistent pooled connections with a call-home handshake and heartbeat-based health monitoring. Backpressure relies on TCP flow control.
  \item The \textbf{parser zoo} contains 17 named parsers mapping to 9 distinct implementations behind enum-dispatched and trait-object interfaces. The proliferation is driven by the absence of a standard tool-call format across model families.
  \item \textbf{Tool call parsing} is a streaming problem where boundary misclassification leads to data loss or corruption---the dominant bug category found by fuzzing. Of the 16 bugs cataloged in Chapter~\ref{ch:findings}, the majority reside in parsers.
\end{itemize}

Table~\ref{tab:parser-bugs} summarizes the parser bugs discussed in this chapter, organized by severity.

\medskip
\begin{table}[h]
\centering
\small
\begin{tabular}{llll}
\hline
\textbf{Bug} & \textbf{Parser} & \textbf{Category} & \textbf{Severity} \\
\hline
\#1  & GLM-4.7           & UTF-8 boundary panic       & Crash \\
\#13 & XML strip\_quotes  & Slice range panic          & Crash \\
\#3  & TwoPartCodec      & Integer overflow           & Critical \\
\#8  & Granite           & Streaming mismatch         & Correctness \\
\#9  & MiniMax           & Streaming mismatch         & Correctness \\
\#11 & Kimi K2           & OnceLock stale cache       & Logic \\
\#10 & DSML              & Silent parameter drop      & Logic \\
\#7  & Pythonic          & Prefix leak                & Logic \\
\#12 & XML literal\_eval  & String value corruption    & Correctness \\
\hline
\end{tabular}
\caption{Parser and codec bugs found by fuzzing, discussed in this chapter.}
\label{tab:parser-bugs}
\end{table}

\medskip
\begin{center}\rule{0.5\linewidth}{0.4pt}\end{center}
\medskip

The codec and parser layers occupy opposite ends of the complexity spectrum: the codec is a single, well-specified state machine with one critical bug; the parser zoo is a sprawling collection of ad-hoc text processors with bugs in nearly every implementation. Both are essential to Dynamo's operation, and both benefit enormously from fuzzing.

We have now mapped every component of Dynamo's architecture, from the transformer inside the GPU through the distributed system that orchestrates a fleet. The next two chapters shift from understanding the system to breaking it: Chapter~\ref{ch:fuzzing} introduces the fuzzing methodology, and Chapter~\ref{ch:findings} catalogs what we found.

\medskip
\noindent\textbf{Check Your Understanding}

\smallskip
\noindent\emph{These questions are answerable from this chapter alone.}

\begin{enumerate}
  \item Why does Dynamo use a custom binary protocol rather than HTTP/JSON or Protocol Buffers for inter-service communication? What latency characteristics make this choice necessary?

  \emph{Answer: During autoregressive decoding, Dynamo generates millions of messages per second (one token per request per step, at thousands of concurrent requests). JSON parsing imposes per-field overhead; Protocol Buffers require schema management and code generation. Dynamo's length-prefixed binary format allows the receiver to determine message boundaries from a fixed 24-byte preamble without parsing any content, enabling zero-allocation reads and zero-copy buffer sharing via the \texttt{bytes} crate.}

  \item Explain how the TwoPartCodec's two-phase design (preamble then payload) enables efficient reading from a TCP stream.

  \emph{Answer: In phase 1, the codec reads exactly 24 bytes to learn the header and body lengths. If fewer than 24 bytes are available, it returns \texttt{Ok(None)} and Tokio waits for more data. In phase 2, it reads exactly \texttt{header\_len + body\_len} bytes. This two-phase approach means the codec never reads more data than needed, never speculatively parses partial messages, and can pre-allocate the exact buffer size---avoiding both over-allocation and reallocation.}

  \item What is the integer overflow vulnerability in the codec's preamble handling? Why does it manifest differently in debug vs.\ release builds?

  \emph{Answer: The computation \texttt{24 + header\_len + body\_len} can overflow when \texttt{header\_len} and \texttt{body\_len} are attacker-controlled \texttt{u64} values cast to \texttt{usize}. In debug builds, Rust's default arithmetic panics on overflow, causing a denial-of-service crash. In release builds, Rust uses wrapping (two's complement) arithmetic by default, so the overflow produces a small \texttt{total\_len} that passes the \texttt{max\_message\_size} check, allowing the attacker to bypass length validation.}

  \item Why is tool call parsing particularly prone to bugs? What makes the streaming (token-by-token) constraint harder than batch parsing?

  \emph{Answer: In batch mode, the parser sees the complete input and can use global search to find delimiters. In streaming mode, it must classify each token as it arrives---inside a tool call, outside, or at a boundary---without future context. Delimiters that span token boundaries may be detected at different points than in batch mode. A single misclassified token causes synchronization loss: the parser either buffers normal text (increasing latency) or emits tool call content as raw text (causing data corruption).}

  \item The \texttt{try\_literal\_eval} function in the XML parser performs \texttt{.replace("True", "true")}. Why is this problematic, and what is a correct alternative?

  \emph{Answer: The replacement is global and not context-aware---it corrupts substrings within quoted string values (e.g., ``TrueNorth'' becomes ``trueNorth''). A correct implementation would only replace Python keywords that appear as standalone tokens outside of string literals. This could be done by first tokenizing the input to identify string boundaries, then only replacing bare keywords, or by using a proper Python literal parser.}

  \item How does the GLM-4.7 parser's function name extraction lead to a UTF-8 boundary panic? What is the root cause?

  \emph{Answer: The parser extracts the function name by trimming whitespace from the content between \texttt{<tool\_call>} and \texttt{<arg\_key>}. It then uses the trimmed name's byte length as an offset into the untrimmed content: \texttt{content[function\_name.len()..]}. When there is leading whitespace and the function name contains multi-byte UTF-8 characters, the byte offset is wrong (off by the whitespace bytes), and may land in the middle of a multi-byte code point. Rust's string slicing requires char-aligned boundaries and panics otherwise.}

  \item Explain the \texttt{OnceLock} caching bug in the Kimi K2 parser. Under what deployment scenario does it cause failures?

  \emph{Answer: The parser stores its compiled regex in a \texttt{static OnceLock}, which initializes exactly once per process. The regex pattern is built from the parser's configuration tokens. If a single process serves multiple Kimi K2 model variants with different delimiter configurations, the regex cached from the first variant is used for all subsequent variants. Tool calls from variants with different delimiters fail to match the cached regex, causing silent tool call detection failures.}

  \item Why does the Pythonic parser use the \texttt{rustpython-parser} crate to parse tool calls, and what are its limitations?

  \emph{Answer: The Pythonic format (\texttt{[func(key=val)]}) is valid Python syntax, so using a real Python parser avoids reimplementing Python's expression grammar---it correctly handles nested lists, dicts, keyword arguments, and Python literals (True, False, None) natively. The limitations are: (1) the regex that detects tool calls requires at least one argument, rejecting zero-argument calls like \texttt{[get\_time()]}; (2) the normal-text extraction uses \texttt{split().next()}, which only captures text before the first match, silently dropping text after the tool call.}
\end{enumerate}
